\chapter{Planning as an Architectural Control}
\label{ch:planning}

\section{Why This Chapter Must Be More Honest}

Planning is central to the manuscript, but the chapter must not pretend to have
proved more than the framework can support. Chapter~\ref{ch:formal-language}
already established the book's honesty policy: we are working with a
qualitative barrier scaffold, not a calibrated behavioural law. The present
chapter therefore proves only those claims that follow from the chapter's
definitions and monotone assumptions. Broader practical claims remain
interpretive and are marked as such.

\section{Definitions and Assumptions}

\begin{definition}[Plan]
A \textbf{plan} for transition $s_0 \to s_1$ is a pre-action specification of
at least one future decision window, one entry action, and one relevant
environmental arrangement for the target state.
\end{definition}

\begin{definition}[Resolved entry specification]
Let $U(s_1)$ denote the target ambiguity from
Chapter~\ref{ch:formal-language}. A plan is said to provide a
\textbf{resolved entry specification} if it fixes a visible first action for
entering $s_1$ and thereby lowers the unresolved specification burden.
\end{definition}

\begin{definition}[Execution choice complexity]
For a given future decision window, let $D$ denote the number of
action-relevant micro-decisions still left to be made at the moment of action.
Examples include choosing the location, selecting the first task, locating the
materials, or deciding when to stop.
\end{definition}

\begin{remark}
The variable $D$ is introduced as a qualitative bookkeeping device, not a
measured cognitive constant. The only property needed below is monotonicity:
if more in-the-moment decisions remain unresolved, then the execution burden is
harder rather than easier.
\end{remark}

\begin{axiom}[Monotone planning assumption]
Holding the current state fixed, the activation barrier
$\DV(s_0 \to s_1)$ is weakly increasing in target ambiguity and weakly
decreasing in relevant preparation.
\end{axiom}

\begin{axiom}[Decision-load assumption]
Holding other terms fixed, the in-the-moment execution burden is weakly
increasing in $D$.
\end{axiom}

\begin{discussion}
These assumptions do not prove that a specific numerical planner outperforms a
specific spontaneous agent in every situation. They only formalize the business
logic already used by the book: leaving fewer questions unanswered makes future
entry easier, and resolving relevant setup in advance lowers the burden at the
moment of action.
\end{discussion}

\section{Planning Before the Moment of Action}

\begin{proposition}[Resolved entry lowers ambiguity]
If a plan provides a resolved entry specification for the target state $s_1$,
then the target ambiguity term $U(s_1)$ is weakly lower than under an otherwise
similar unplanned intention.
\end{proposition}

\begin{proof}
By definition, a resolved entry specification fixes a visible first action that
would otherwise remain undecided at the moment of action. Since $U(s_1)$ was
defined as the unresolved specification burden of entering $s_1$, resolving at
least one such burden weakly lowers $U(s_1)$.
\end{proof}

\begin{proposition}[Planning lowers future barrier under the monotone scaffold]
Suppose a plan for $s_0 \to s_1$ both lowers target ambiguity and increases
relevant preparation. Under the monotone planning assumption, the future
activation barrier is weakly lower than without that plan.
\end{proposition}

\begin{proof}
From the previous proposition, resolved entry lowers $U(s_1)$. By definition
of planning, relevant preparation is also increased. The monotone planning
assumption states that $\DV(s_0 \to s_1)$ is weakly increasing in ambiguity and
weakly decreasing in preparation. Therefore lowering ambiguity and increasing
preparation weakly lowers the activation barrier.
\end{proof}

\begin{corollary}[Last-minute self-persuasion is structurally weaker]
If last-minute self-persuasion changes neither ambiguity nor preparation, while
planning changes at least one of them in the favourable direction, then the
planned transition has weakly lower activation barrier than the merely
persuaded transition.
\end{corollary}

\begin{proof}
The merely persuaded transition leaves the relevant barrier variables unchanged
by assumption. The planned transition lowers at least one favourable term.
Apply the previous proposition.
\end{proof}

\begin{remark}
This corollary is deliberately weaker than the chapter's earlier ``planning
theorem.'' It does not claim that planning always wins in every human setting.
It states only that if planning actually changes the variables that the model
says matter, then the model predicts a lower barrier than verbal endorsement
alone.
\end{remark}

\section{A Planning Stack}

For teaching purposes, it is useful to separate planning into layers:
\begin{enumerate}[nosep]
  \item \textbf{Outcome layer:} what end state is desired?
  \item \textbf{Session layer:} what concrete session will move that outcome?
  \item \textbf{Entry layer:} what exact first action starts the session?
  \item \textbf{Environment layer:} what physical arrangement lowers entry
        friction?
  \item \textbf{Trigger layer:} when and where is the decision window expected?
\end{enumerate}

\begin{proposition}[Execution-choice compression]
If a plan resolves $k \geq 1$ future micro-decisions before the decision
window, then the execution choice complexity at the window is reduced from $D$
to at most $D-k$.
\end{proposition}

\begin{proof}
By construction, each resolved micro-decision is one less action-relevant item
to be decided at execution time. Therefore resolving $k$ such items before the
window reduces the remaining decision count by at least $k$, so the future
decision complexity is at most $D-k$.
\end{proof}

\begin{corollary}[Binary execution advantage, qualified]
If a plan reduces execution complexity to a binary choice---execute the
predefined action or refuse it---then the in-the-moment execution burden is
weakly lower than under any otherwise similar plan with strictly larger
execution complexity.
\end{corollary}

\begin{proof}
A binary choice has execution complexity $D=1$. Any otherwise similar plan with
strictly larger complexity has $D>1$. By the decision-load assumption, greater
execution complexity weakly increases in-the-moment burden. Therefore the
binary-choice plan weakly lowers that burden.
\end{proof}

\begin{discussion}
This result is narrower and more honest than the older formulation. It does not
claim that binary framing guarantees compliance. It claims only that, within
the model, reducing the number of unresolved real-time decisions reduces the
decision burden that must be paid at the window.
\end{discussion}

\section{Minimal Planning Workflow}

\begin{templatebox}
\textbf{Minimal transition plan}
\begin{enumerate}[nosep]
  \item Name the target state.
  \item Name the likely decision window.
  \item Write the first visible action.
  \item Place the required objects in the future environment.
  \item Define the stop rule for the session.
\end{enumerate}
\end{templatebox}

\begin{remark}
The workflow is pedagogical rather than theorem-like. Its justification is that
each line targets one of the variables formalized above: target state fixes the
objective, decision window fixes timing, first action lowers ambiguity,
environment placement raises preparation, and stop rule limits in-session
redesign cost.
\end{remark}

\section{Why Unplanned Intentions Often Fail}

\begin{proposition}[Unplanned intention preserves current advantage]
Assume two future options compete at a decision window: remaining in the
present state $s_0$ or attempting transition to $s_1$. If no plan lowers the
barrier terms for $s_1$, then the model provides no new structural reason for
$s_1$ to become easier relative to the status quo at that window.
\end{proposition}

\begin{proof}
By assumption, no plan lowers ambiguity, raises preparation, or otherwise
reduces the transition burden for $s_1$. Therefore the model's barrier-relevant
variables are unchanged in favour of $s_1$. Since the current state's
structural advantage has not been reduced by planning, the model introduces no
new formal reason for the transition to become easier relative to the status
quo.
\end{proof}

\begin{corollary}[Unplanned preference alone does not reverse inertia]
Within this framework, a merely verbal preference for $s_1$ does not by itself
imply that $s_1$ will defeat the current state's structural advantage at the
next decision window.
\end{corollary}

\begin{proof}
Apply the previous proposition to the special case where the only new element
is verbal endorsement, and no barrier-relevant planning change has occurred.
\end{proof}

\begin{example}[Exercise after work]
The sentence ``I should work out after work'' is not yet a usable plan. It
leaves open the location, bag, sequence, exercise selection, session length,
and exit condition. In the chapter's notation, these unresolved items keep
execution complexity high and leave target ambiguity largely unchanged. A real
plan resolves some of those variables upstream.
\end{example}

\section{Evidence and Reading Discipline}

This chapter remains primarily a modelling chapter. Its formal claims are
proved only relative to:
\begin{enumerate}[nosep]
  \item the chapter's definitions;
  \item the monotone barrier scaffold from Chapter~\ref{ch:formal-language};
  \item the decision-load assumption introduced here.
\end{enumerate}

It does \emph{not} claim laboratory proof that every human plan reduces action
cost in every domain. The chapter proves what follows inside the model and no
more.

\section{Recommended Reading}

\begin{readingbox}
\textbf{Foundation}
\begin{itemize}[nosep]
  \item Steven Strogatz, \emph{Nonlinear Dynamics and Chaos}, for the logic of
        regime change and why small structural changes can matter more than raw
        force.
  \item Eugene Izhikevich, \emph{Dynamical Systems in Neuroscience}, for a more
        disciplined state-transition vocabulary than generic self-management
        language provides.
\end{itemize}
\textbf{Model discipline}
\begin{itemize}[nosep]
  \item Revisit Chapter~\ref{ch:formal-language} before reading this chapter as
        if it had established a calibrated behavioural law. It has not.
  \item Read Chapter~\ref{ch:dynamics} together with this chapter: planning is
        valuable because it changes later barrier structure before the decision
        window opens.
\end{itemize}
\textbf{Boundary}
\begin{itemize}[nosep]
  \item Empirical consciousness sources motivate taking structured state
        transitions seriously, but they are not direct proof of this chapter's
        planning propositions.
  \item Frontier physical theories are not needed to justify planning as an
        architectural control.
\end{itemize}
\end{readingbox}

\begin{summarybox}
This chapter now does less pretending and more actual work. It defines what a
plan changes, proves that resolving ambiguity and increasing preparation lower
future barrier under the model's monotone assumptions, proves a qualified
binary-execution advantage, and downgrades broader behavioural claims to the
level of disciplined interpretation. That is formally narrower than the
chapter's earlier rhetoric, but it is much closer to textbook honesty.
\end{summarybox}
